% Revised 2026-10-17: conditional evidence return; changes re-enter protocol.
% export-reference-width-mm: 169
\begin{tikzpicture}[x=1mm,y=1mm,
  font=\figurefont\fontsize{8}{10}\selectfont,
  box/.style={rounded corners=1.8pt,draw=FigureStroke,fill=FigureBlueFill,
    line width=.6pt,text width=29mm,minimum height=15mm,
    align=center,inner sep=2mm},
  flow/.style={-{Latex[length=2mm]},draw=FigureStroke,line width=.75pt},
  lab/.style={font=\FigureNoteFont\fontsize{7.5}{9.5}\selectfont,
    text=FigureInk,inner sep=.6mm}]
  \node[box,fill=FigureYellowFill] (protocol) at (17,0)
    {预设要求与协议\\对象、分布、阈值};
  \node[box] (evidence) at (61,0) {取得证据\\评价记录、模型依据};
  \node[box] (bound) at (105,0) {计算指标与误差界\\检查适用前提};
  \node[box] (judge) at (149,0) {达标判断\\通过／否决／不足};
  \node[box,text width=43mm,fill=FigureYellowFill] (change) at (39,-38)
    {策略、环境或保护机制改变\\重新检查原协议是否适用};
  \node[box,text width=43mm] (act) at (127,-38)
    {补采评价回合、收窄范围\\保护或回退};
  \draw[flow] (protocol.east) -- (evidence.west);
  \draw[flow] (evidence.east) -- (bound.west);
  \draw[flow] (bound.east) -- (judge.west);
  \draw[flow] (judge.south) -- (149,-19) -- (127,-19) -- (act.north);
  % Additional evidence may reuse an unchanged protocol; changed objects
  % or protocol must be checked first, never bypassing the protocol node.
  \draw[flow,dashed] (act.north west) -- (103,-17)
    -- (61,-17) -- (evidence.south);
  \node[lab,above] at (82,-17) {原协议不变时补采};
  \draw[flow] (act.west) -- node[lab,above] {对象或协议改变} (change.east);
  \draw[flow] (change.north) -- (39,-20) -- (17,-20) -- (protocol.south);
\end{tikzpicture}
